\section{Methodology}
\label{sec:methodology}

We compare three prompting conditions applied to the same model (\gptfour via the OpenAI API) across three evaluation tasks.
All conditions share the same model, decoding parameters, and token budgets; only the system prompt differs.

\subsection{Prompting Conditions}
\label{sec:conditions}

\para{\standard.}
The system prompt is simply ``You are a helpful assistant.'' No special reasoning instructions are provided.

\para{\standardcot.}
The system prompt instructs the model to ``think step by step before giving your answer.'' This follows the conventional CoT paradigm~\citep{wei2022chain}, in which all reasoning appears before the final response.

\para{\iscotshort (ours).}
The system prompt requires the model to generate a \texttt{<think>} block before \emph{every sentence} of its visible response.
Each \texttt{<think>} block must consider (1) the accuracy of the upcoming sentence, (2) its phrasing, and (3) its logical connection to the preceding sentence.
The model alternates between \texttt{<think>[reasoning]</think>} and visible text, producing a pattern in which every claim is preceded by explicit deliberation.
We strip the \texttt{<think>} blocks during evaluation so that only the visible text is assessed.
This design leverages the model's native instruction-following and CoT capabilities without any training modifications.

\subsection{Evaluation Tasks}
\label{sec:tasks}

\para{Open-ended generation (Experiment~1).}
We curate 30 diverse prompts covering explanations, essays, and advice (\eg ``Explain why the sky appears blue during the day but red at sunset'').
Each prompt is answered under all three conditions, yielding 90 responses.
We evaluate using (a) automated text analysis measuring hedging, self-correction, lexical diversity, and sentence structure; and (b) LLM-as-judge pairwise comparisons on coherence, depth, and accuracy, with response order randomized to control for positional bias.

\para{Factual accuracy (Experiment~2).}
We sample 100 questions from \truthfulqa~\citep{lin2022truthfulqa} (seed = 42, from 817 available).
Each question is answered under all three conditions, yielding 300 responses.
We evaluate truthfulness and informativeness using an LLM judge following the \truthfulqa evaluation protocol.

\para{Math reasoning (Experiment~3).}
We sample 100 problems from the \gsmk test set~\citep{cobbe2021gsm8k} (seed = 42, from 1{,}319 available).
Each problem is solved under all three conditions, yielding 300 responses.
We evaluate using exact match on the final numerical answer.

\para{Cross-experiment text analysis (Experiment~4).}
We aggregate all 130 prompts per condition across Experiments 1--3 and compute text-level properties: hedging rate, self-correction rate, qualifying clause rate, type-token ratio (\ttr), and average sentence length.

\subsection{Hyperparameters}
\label{sec:hyperparameters}

\begin{table}[t]
\centering
\caption{Hyperparameter settings across experiments. Temperature is set to 0.0 for tasks requiring deterministic output (math, judging) and 0.7 for open-ended generation.}
\begin{tabular}{@{}lll@{}}
\toprule
\textbf{Parameter} & \textbf{Value} & \textbf{Scope} \\
\midrule
Model & \gptfour & All \\
Temperature & 0.7 & Open-ended \\
Temperature & 0.0 & \gsmk, judge \\
Max tokens & 1{,}024 & Open-ended, \gsmk \\
Max tokens & 512 & \truthfulqa \\
Random seed & 42 & All \\
\bottomrule
\end{tabular}
\label{tab:hyperparameters}
\end{table}

\subsection{Statistical Tests}
\label{sec:stats}

For continuous text metrics, we use the Wilcoxon signed-rank test (paired, non-parametric) and report Cohen's $d$ for effect size.
For binary accuracy outcomes (\truthfulqa, \gsmk), we use McNemar's test over discordant pairs.
All tests use a significance level of $\alpha = 0.05$.

\subsection{Text Analysis Metrics}
\label{sec:text_metrics}

We define four categories of sentence-level linguistic markers, detected using keyword lists applied via NLTK tokenization:

\para{Hedging language.}
Words and phrases signaling uncertainty or tentativeness (\eg ``might,'' ``perhaps,'' ``it is possible that'').
We report the rate per sentence.

\para{Self-correction markers.}
Phrases indicating the model is revising or qualifying a prior statement (\eg ``actually,'' ``however,'' ``on second thought'').

\para{Qualifying clauses.}
Subordinate clauses that narrow the scope of a claim (\eg ``in most cases,'' ``under certain conditions'').

\para{Lexical diversity.}
Type-token ratio (\ttr), defined as the number of unique word types divided by the total number of word tokens.
Higher \ttr indicates greater vocabulary variety.
